% mono/preamble/environments.tex
%
% Segment environments — the visual register for the 19 segment types
% defined in FORMAT.md. Architecture:
%
%   \segmenthead{<Type Label>}{<Title>}{<status>}
%      ... segment body ...
%   \segmentfoot
%
% The kramdown emitter produces these directly. No 19 distinct LaTeX
% environments — one parameterized macro pair, type label as a string. This
% keeps the AAD type taxonomy and the LaTeX taxonomy decoupled: when FORMAT
% grows a 20th type, the build pipeline maps it to the right label and the
% rendering needs no change here.
%
% Numbering uses a single shared counter per chapter (matches the existing
% bin/build markdown behavior: "Definition I.3", "Result II.7", etc.). The
% counter is named 'segment' rather than 'claim' (per FORMAT's vocabulary —
% segments hold claims; not all segments contain a single claim).

\newcounter{segment}[chapter]
% FORMAT.md uses "Definition I.3, Result II.7" with roman chapter numerals;
% match it here. (kaobook's kao chapter style already renders the BIG margin
% chapter glyph in roman, but \thechapter itself is arabic — \Roman{chapter}
% in the counter format aligns the body convention with the head.)
\renewcommand{\thesegment}{\Roman{chapter}.\arabic{segment}}

% Internal section subheads — \segmentsubhead{Formal Expression}, etc.
\NewDocumentCommand{\segmentsubhead}{m}{%
  \par\medskip%
  \noindent{\sffamily\bfseries #1}\par\nobreak\smallskip%
}

% Segment opener. Three args:
%   #1 = type label (string, e.g. "Result", "Derivation")
%   #2 = title (the human-readable form of the slug)
%   #3 = status (string, e.g. "exact"; empty string = no badge)
%
% The third argument is always required; the kramdown emitter passes the
% empty string when frontmatter has no `status:` set, which suppresses the
% badge via the IfStrEq below.
\NewDocumentCommand{\segmenthead}{m m m}{%
  \refstepcounter{segment}%
  \par\medskip%
  \noindent\textcolor{rule-gray}{\rule{\linewidth}{0.3pt}}\par%
  \nobreak\smallskip%
  \noindent%
  {\sffamily\bfseries #1~\thesegment}%
  \IfStrEq{#3}{}{}{\statusbadge{#3}}%
  \nobreak\quad\textcolor{rule-gray}{\textbullet}\quad%
  {\itshape #2}.%
  \par\nobreak\smallskip%
  \noindent\textcolor{rule-gray}{\rule{\linewidth}{0.3pt}}\par%
  \nobreak\medskip%
}

% Segment closer.
\NewDocumentCommand{\segmentfoot}{}{%
  \par\medskip%
  \noindent\textcolor{rule-gray}{\rule{0.3\linewidth}{0.3pt}}\par%
  \medskip%
}

% Stage marker — review-mode only. Surfaces the `stage:` frontmatter field
% (draft / deps-verified / claims-verified / format-clean / candidate) as a
% small marginal note. The kramdown emitter only calls this in :review
% variant; :public elides it entirely.
\NewDocumentCommand{\segmentstage}{m}{%
  \marginnote{\footnotesize\sffamily\color{status-tentative}stage: #1}%
}

% Missing-segment placeholder — for slugs listed in OUTLINE.md but without
% a corresponding src/ file. Generates a TOC entry and a visible marker so
% reviewers see the gap in the rendered output, not just in a sidecar list.
\NewDocumentCommand{\missingsegment}{m m m}{%
  % #1 = type label, #2 = slug, #3 = claim text (from the outline table)
  \refstepcounter{segment}%
  \par\medskip%
  \noindent\textcolor{rule-gray}{\rule{\linewidth}{0.3pt}}\par%
  \nobreak\smallskip%
  \noindent%
  {\sffamily\bfseries\color{status-tentative} #1~\thesegment}%
  \quad\textcolor{rule-gray}{\textbullet}\quad%
  {\itshape\color{status-tentative} #3}.%
  \par\nobreak\smallskip%
  \noindent\emph{Segment \texttt{#2} not yet written.}\par%
  \nobreak\smallskip%
  \noindent\textcolor{rule-gray}{\rule{\linewidth}{0.3pt}}\par%
  \nobreak\medskip%
}

% Gap marker — for --GAP-- entries in OUTLINE tables (theory not yet
% developed; no slug, just a description of the open question).
\NewDocumentCommand{\outlinegap}{m}{%
  \par\medskip%
  \noindent\textcolor{status-tentative}{\textbf{[Gap]}}~\emph{#1}\par%
  \medskip%
}
